\documentclass[11pt]{article}
\usepackage[margin=1in]{geometry}
\usepackage{amsmath, amssymb, amsthm, bm, mathtools}
\usepackage{booktabs, array}
\usepackage{enumitem}
\usepackage{siunitx}
\usepackage{hyperref}
\hypersetup{colorlinks=true, linkcolor=blue, urlcolor=blue, citecolor=blue}
\sisetup{round-mode=places,round-precision=2,group-separator={,}}

\newcommand{\dd}{\,\mathrm{d}}
\newcommand{\R}{\mathbb{R}}

\title{Practice Problem Set 1}
\author{ECON 308 -- International Economic Relations}
\date{Fall 2025}

\begin{document}
\maketitle



\begin{enumerate}[leftmargin=*]
\item Explain why country size (GDP) and geographic distance are strong predictors of bilateral trade flows. Provide at least one microeconomic channel for each.
\item Define ``home bias'' in trade. Describe one empirical method for detecting a border effect in a gravity regression and one interpretation of finding a large border effect.
\item Why do small economies often exhibit higher trade-to-GDP ratios than large economies? Include at least one scale or variety argument.
\item Chapter 2 highlights both inter-industry and intra-industry trade. Briefly differentiate them and give one example of each between advanced economies.
\end{enumerate}


\begin{enumerate}[leftmargin=*, resume]
\item Define both. Explain, using unit labor requirements, how a country can gain from trade even if it has an absolute disadvantage in \emph{both} goods.
\item In the Ricardian model with free trade, what pins down relative wages across countries? Explain the role of world prices and specialization.
\item Distinguish between the production possibility frontier (PPF) and the consumption possibility frontier (CPF) under autarky versus free trade. Why does the CPF typically lie outside the PPF with trade?
\item Briefly explain why, under perfect competition, autarky relative prices equal unit labor cost ratios, and how opening to trade changes the mapping from technology to prices and wages.
\end{enumerate}
\bigskip



\begin{enumerate}[leftmargin=*, resume]
\item
The (log) gravity model is
\[
\ln T_{ij}=\alpha + \beta_1 \ln Y_i + \beta_2 \ln Y_j - \gamma \ln D_{ij},
\]
with $T_{ij}$ bilateral exports from $i$ to $j$, $Y$ GDP, and $D$ great-circle distance (km). Suppose $\beta_1=\beta_2=1$, $\gamma=1$, and $e^\alpha=10^{-6}$ so that
\[
T_{ij}=10^{-6}\,\frac{Y_i Y_j}{D_{ij}}.
\]
Consider the following economies (values in billions of US dollars; distances in km):
\[
\begin{array}{@{}lrrrr@{}}
\toprule
\text{Country} & Y \ (\$\,\text{bn}) & & & \\
\midrule
A & 2{,}000 &&&\\
B & 1{,}000 &&&\\
C & \phantom{0}500 &&&\\
\bottomrule
\end{array}
\qquad
\begin{array}{@{}lrrr@{}}
\toprule
 & A & B & C\\
\midrule
A & - & 6{,}000 & 1{,}500\\
B & 6{,}000 & - & 2{,}000\\
C & 1{,}500 & 2{,}000 & - \\
\bottomrule
\end{array}
\]
\begin{enumerate}
\item Compute the predicted bilateral exports $T_{AB}$, $T_{AC}$, and $T_{BC}$.
\item By what \emph{percent} does $T_{AB}$ change if $D_{AB}$ falls by 10\%? (Use the elasticity implied by the model.)
\item If both $Y_A$ and $Y_B$ rise by 5\%, what is the approximate percentage change in $T_{AB}$?
\item Suppose the true $T_{AB}$ observed in data is \$\num{1.3} billion. Compute the (log) residual $\ln T^{\text{obs}}_{AB}-\ln T^{\text{pred}}_{AB}$ and provide one economic interpretation of a positive residual.
\end{enumerate}

\item
You are given the following national accounts (all in billions of domestic currency) and nominal GDP:
\[
\begin{array}{@{}lrrrr@{}}
\toprule
\text{Country} & \text{Exports} & \text{Imports} & \text{GDP} & \text{Population (mn)}\\
\midrule
D & 120 & 100 & 1{,}000 & 50\\
E & 400 & 350 & 2{,}500 & 200\\
F & 70 & 110 & \phantom{0}600 & 10\\
\bottomrule
\end{array}
\]
\begin{enumerate}
\item Compute each country's openness ratio $(X+M)/\text{GDP}$ (in \%).
\item Compute \emph{per capita} exports and imports for each country.
\item Rank countries by openness and briefly explain why small economies often have higher openness ratios than large ones.
\end{enumerate}

\item 
Extend the gravity model in Q1 by multiplying predicted trade by $(1+\lambda_{CL})$ if a common language is shared and by $(1-\lambda_{BOR})$ if there is an international border friction.\footnote{Here, ``border friction'' means $i$ and $j$ are in different economic unions with additional administrative costs. Choose the signs so that a border \emph{reduces} trade.}
Let $\lambda_{CL}=0.20$ and $\lambda_{BOR}=0.30$. Countries $A$ and $B$ share a common language; $A$ and $C$ do not; all pairs are separated by an international border.
\begin{enumerate}
\item Starting from your baseline $T_{ij}$ in Q1(a), compute adjusted $T_{AB}$, $T_{AC}$, and $T_{BC}$.
\item By what percent does the common language increase $T_{AB}$ relative to the case with no language effect?
\item Holding other parameters fixed, solve for the value of $\lambda_{BOR}$ that would exactly offset the common language boost between $A$ and $B$.
\end{enumerate}

\item 
Suppose over a decade countries $A$ and $C$ experience average annual real GDP growth rates of 2.5\% and 4\%, respectively, while $D_{AC}$ is unchanged. Assume prices are stable and growth maps one-for-one into $Y$.
\begin{enumerate}
\item Compute $Y_A$ and $Y_C$ after 10 years.
\item Using the gravity model with $\beta_1=\beta_2=1$, $\gamma=1$, compute the growth factor for $T_{AC}$ over the decade.
\item If observed $T_{AC}$ grows by only 12\%, list two trade-cost or policy changes consistent with the shortfall relative to the gravity prediction.
\end{enumerate}
\end{enumerate}







\begin{enumerate}[leftmargin=*, resume]
\item 
Two countries, Home ($H$) and Foreign ($F$), produce Wine ($W$) and Cloth ($C$) with one input, labor. Unit labor requirements are:
\[
\begin{array}{@{}lcc@{}}
\toprule
 & a_{LW} & a_{LC} \\
\midrule
H & 2 & 4 \\
F & 6 & 3 \\
\bottomrule
\end{array}
\qquad
\text{Labor endowments: } L_H=240,\; L_F=240.
\]
\begin{enumerate}
\item Compute each country’s autarky opportunity cost of $W$ in terms of $C$ and vice versa. Identify comparative advantage.
\item Draw (or describe algebraically) each country’s PPF. What is each country’s autarky relative price $P_W/P_C$ under competitive pricing equal to unit labor cost?
\item Find the range of world relative prices $P_W/P_C$ consistent with complete specialization.
\item Suppose the world relative price is $P_W/P_C=1.2$. Determine specialization, production bundles, and the pattern of trade.
\item Compute the real wage in each country under free trade in terms of each good (units of $W$ and units of $C$ purchasable per unit of labor income).
\end{enumerate}

\item 
Continue from Q1 with $P_W/P_C$ moving from $1.2$ to $0.8$.
\begin{enumerate}
\item Re-evaluate which country (if any) changes specialization as the relative price falls to $0.8$. Explain.
\item Compute the real wage of each country in terms of $W$ and $C$ at $P_W/P_C=0.8$ and compare to autarky.
\item Which country’s real wage benefits from an improvement in its terms of trade? Explain using the Ricardian logic.
\end{enumerate}

\item 
There are three countries with the following unit labor requirements:
\[
\begin{array}{@{}lcc@{}}
\toprule
 & a_{LW} & a_{LC} \\
\midrule
1 & 1 & 2 \\
2 & 2 & 3 \\
3 & 4 & 5 \\
\bottomrule
\end{array}
\quad L_1=L_2=L_3=100.
\]
\begin{enumerate}
\item Order countries by comparative advantage in $W$ relative to $C$.
\item Derive the world relative supply (WRS) schedule $RS^{W/C}(P_W/P_C)$ by progressively ``turning on'' countries into $W$ as $P_W/P_C$ rises. State the price thresholds where kinks occur.
\item If world relative demand is $RD^{W/C}=1.1 - 0.2(P_W/P_C)$ (for $P_W/P_C\in[0,5]$), solve for the free-trade equilibrium $P_W/P_C$ and identify which countries produce which good.
\end{enumerate}

\item
In the two-country model of Q1, suppose Home experiences a neutral productivity improvement in Cloth only: $a_{LC}$ falls from $4$ to $3$ (all else unchanged).
\begin{enumerate}
\item Re-compute Home’s opportunity cost of $W$ in terms of $C$ and the specialization range.
\item Discuss how the shock shifts the world relative supply and the likely effect on the world relative price $P_W/P_C$.
\item With $P_W/P_C$ adjusting to the new equilibrium, compare Home and Foreign real wages (in both goods) to the pre-shock free-trade case. Who gains more from this productivity growth?
\end{enumerate}
\end{enumerate}



\end{document}
